\section{Formation as constrained optimal control}\label{sec:control}

\subsection{State, control and a dimensionally defined objective}
Let $\bm\vartheta$ collect fixed but possibly uncertain material and apparatus
parameters. The symbol is distinct from temperature $\Theta$ and contact angle
$\theta$. Let $\mathcal D$ denote the full dynamic residual, including bulk
equations, moving-interface laws, boundary conditions and algebraic constraints.
For a prescribed initial physical state $\bm y_0$, write
\begin{equation}
 \mathcal D(\bm y,\dot{\bm y},\bm a;\bm\vartheta)=0,
 \qquad \bm y(0)=\bm y_0,\qquad 0\leq t\leq T_f.
 \label{eq:dynamic-constraint}
\end{equation}
This notation can represent the full compressible problem or a declared
averaged reduction. In the former the command is the actual waveform; in the
latter it may be its slowly varying envelope. They are not interchangeable
for an arbitrarily short pulse.

For shapes close enough to the target to admit a normal correspondence, let
$\eta_s$ be their signed normal departure on $\Gamma_{\rm opt,\star}$. Define
the target optical area $A_{\rm opt,\star}$ and a positive geometric tolerance
$h_{\rm ref}$. A dimensionless optional shape error is
\begin{equation}
 J_s=\frac{1}{A_{\rm opt,\star}}
       \int_{\Gamma_{\rm opt,\star}}\frac{\eta_s^2}{h_{\rm ref}^2}\dd A.
 \label{eq:shape-cost}
\end{equation}
For large departures use a specified geometric distance with an explicit
correspondence or level-set convention; there is no implicit registration that
moves the conjugates. The optical objective $J_{\rm opt}$ is a declared
nonnegative combination of $J_W,J_d,J_{\rm spot}$, evaluated with the actual
index fields if they vary. The simple homogeneous-path expression for $J_W$
must then be replaced by optical path integration through those fields.

Let $\bm S_a$ be a diagonal matrix of positive componentwise command scales
and $T_{\rm ref}>0$ a rate scale in seconds. Let
$\beta_o,\beta_s,\beta_a,\beta_{\dot a}\geq0$ be dimensionless weights and
$\Phi_T(\bm y(T_f))$ a dimensionless terminal objective. A concrete family of
problems is
\begin{align}
 \min_{\bm a,\bm y}\quad
 J&=\Phi_T(\bm y(T_f))
    +\int_0^{T_f}\ell(\bm y,\bm a,\dot{\bm a},t)\dd t,
 \label{eq:dynamic-objective}\\
 \ell&=\frac{1}{T_f}\left[
 \beta_oJ_{\rm opt}+\beta_sJ_s
 +\beta_a\norm{\bm S_a^{-1}\bm a}^2
 +\beta_{\dot a}\norm{T_{\rm ref}\bm S_a^{-1}\dot{\bm a}}^2\right],
 \label{eq:running-cost}
\end{align}
subject to \cref{eq:dynamic-constraint} and the physical constraints below.
The norm penalties describe command effort and smoothness; they are not
automatically acoustic power. A calibrated input-energy term can replace or
supplement them. For a velocity-controlled fluid boundary with domain-outward
normal $\bm n_a$ and wall velocity $\bm v_a$, instantaneous mechanical power
entering that fluid is
$\int_{\Gamma_a}(\bm T\bm n_a)\cdot\bm v_a\dd A$. Electrical power requires
the electromechanical source model. Reversible energy returned to a source and
gross consumed power are different quantities.

Terminal conditions depend on the purpose. For capture after a pulse, choose
$\Phi_T$ to penalize optical error, residual fluid kinetic energy and remaining
acoustic energy, each divided by a declared positive energy or optical scale.
Require the command to end at zero and apply the retention condition in
\cref{eq:unforced-retention}. For driven operation, target the complete holding
state instead; its required nonzero acoustic field is not an unwanted residual.
For curing, add conversion, heat release, shrinkage, evolving rheology and the
final optical index before claiming a retained solid lens.

\subsection{Constraints carry physical meaning}
The admissible set must specify the following quantities, according to the
chosen physical model:
\begin{enumerate}
\item conserved masses or supplied volume constraints, support and wetting
conditions, clearance and an admissible non-self-intersecting interface;
\item command amplitude, bandwidth, rate, total energy, source loading and
the range in which its acoustic response remains valid;
\item acceptable temperature and flow, and any pressure, cavitation or damage
limits supported by material and apparatus data;
\item optical path and direction tolerances at the requested conjugates,
transmission without total internal reflection, and the intended aperture,
wavelength and exposure interval;
\item stability, disturbance rejection and actuator headroom during the
specified holding interval.
\end{enumerate}
There is no universal cavitation threshold or acceptable intensity independent
of liquid, nuclei, frequency and exposure. Such a constraint must come from a
specified model or measurement. Similarly, optical quality after curing cannot
be certified with only liquid-state refractive indices.

\subsection{Adjoint equations give a general derivative method}
The inverse method need not be a new heuristic for each target. A common
constrained state equation and a common objective produce the same sensitivity
and adjoint construction; geometry and constitutive data change its operators.
Represent complex controls by their stacked real and imaginary components
before taking the real derivatives below.

For these operator derivatives, fields on moving domains are pulled back to
fixed reference domains with fixed duality pairings. The residuals then include
the coordinate Jacobians, moving-coordinate transport and transformed boundary
conditions. An equivalent geometric derivation must retain those same transport
and measure-variation terms. Differentiating a fixed-domain equation while
silently changing its domain would not give the required shape adjoint.

For the stationary residual $\mathcal R$ and objective $J_{\rm stat}$, introduce
an adjoint multiplier $\bm\lambda$ in the dual residual space and the Lagrangian
$\mathscr L=J_{\rm stat}+\inner{\bm\lambda}{\mathcal R}$. The superscript
$\ast$ on an operator now denotes its adjoint in the stated real duality
pairing, rather than scalar conjugation. Eliminating state variations gives
\begin{equation}
 \mathcal R_{\bm y}^{\ast}\bm\lambda
        =-(J_{\rm stat})_{\bm y},\qquad
 \nabla_{\bm w}j=(J_{\rm stat})_{\bm w}
                       +\mathcal R_{\bm w}^{\ast}\bm\lambda.
 \label{eq:stationary-adjoint}
\end{equation}
Here $j(\bm w)=J_{\rm stat}(\bm y(\bm w),\bm w)$ is the reduced objective.
The derivatives include all state-dependent loads and boundary terms.
The adjoint approach supplies many control derivatives from a state solve and
an adjoint solve per scalar functional \cite{giles2000}. It does not assert
convexity or global optimality.

For dynamics, introduce the time-dependent adjoint $\bm\Lambda(t)$ and use
\begin{equation}
 \mathscr L_T=\Phi_T(\bm y(T_f))+
   \int_0^{T_f}\left[\ell+
        \inner{\bm\Lambda}{\mathcal D(\bm y,\dot{\bm y},\bm a)}\right]\dd t.
 \label{eq:dynamic-lagrangian}
\end{equation}
Integrating the $\delta\dot{\bm y}$ term by parts gives the formal adjoint
equation and terminal condition
\begin{align}
 \ell_{\bm y}+\mathcal D_{\bm y}^{\ast}\bm\Lambda
       -\frac{\dd}{\dd t}
            (\mathcal D_{\dot{\bm y}}^{\ast}\bm\Lambda)&=0,
 \label{eq:dynamic-adjoint}\\
 (\Phi_T)_{\bm y}
       +\mathcal D_{\dot{\bm y}}^{\ast}\bm\Lambda(T_f)&=0.
 \label{eq:adjoint-terminal}
\end{align}
The terminal equation applies to free terminal state components. Prescribed
terminal constraints add multipliers. Spatial integrations by parts provide
adjoint boundary and interface conditions; these depend on the actual physical
boundaries and objective and cannot be universally omitted. Differential--
algebraic constraints require their consistent adjoint treatment or elimination.

For a fixed formation horizon, the interior command gradient is
\begin{equation}
 \bm g_a=\ell_{\bm a}
       -\frac{\dd}{\dd t}\ell_{\dot{\bm a}}
       +\mathcal D_{\bm a}^{\ast}\bm\Lambda.
 \label{eq:control-gradient}
\end{equation}
A free command endpoint has the associated natural boundary term from
$\ell_{\dot{\bm a}}$; a prescribed pulse endpoint has no free endpoint
variation. If $T_f$ is optimized, its transversality condition must also be
included. State and path constraints augment the Lagrangian with their
multipliers and complementarity conditions. For a convex admissible control
set $\mathcal U$ after those constraints are treated, first-order optimality
requires
\begin{equation}
 \int_0^{T_f}\inner{\bm g_a}{\bm v-\bm a}\dd t\geq0
       \quad\text{for every }\bm v\in\mathcal U.
 \label{eq:control-variational-inequality}
\end{equation}
These are necessary local optimality equations under regularity assumptions;
they are not a closed-form excitation or a proof of a global optimum.

\subsection{Robustness uses one realizable policy}
Let $\mathcal V$ be a declared set of uncertain parameter values, and let
$\pi$ be a causal policy mapping available measurement histories to commands.
A robust formulation is
\begin{equation}
 \min_\pi\ \sup_{\bm\vartheta\in\mathcal V}
       J(\bm y_{\bm\vartheta},\pi;\bm\vartheta),
 \quad\text{with the governing and operating constraints for every }
 \bm\vartheta\in\mathcal V.
 \label{eq:robust-design}
\end{equation}
The same policy acts in every case; it cannot choose a different command using
an unobserved true parameter. Parameter estimation can be part of the
observation state. An open-loop pulse is the special case of a policy fixed
before observing the response. This separates a repeatable inverse-design
method from a shape fitted under a single favorable calibration.
